\section{The Epistemological Conflict: Truthfulness vs. Faithfulness}
\label{sec:conflict}
% Length: ~20%
% Focus: Open the black box. Explain the internal battle between the model's memory and the retrieved text.

% Must contain:
% 1. The Paradox: Define the conflict between parametric knowledge and contextual knowledge.
% 2. Probing Analysis: Synthesize how researchers prove this conflict exists through internal observation or testing.

\subsection{Defining the Paradox}
The foundational vulnerability of automated evaluators lies in the epistemological struggle between inherent parametric memory and externally retrieved evidence. Evaluators frequently fail because they conflate objective factuality with contextual adherence. When a language model judge evaluates a generative output, it intuitively relies on its extensive pretraining corpus. If the provided context contradicts this internal repository, the model experiences severe representational dissonance. 

The framework proposed by \citet{fadeeva2026faithfulnessawareuncertaintyquantificationfactchecking} demonstrates that evaluators must delineate these dimensions by applying distinct uncertainty quantification metrics. They argue that truthfulness requires assessing the internal entropy of the language model, whereas faithfulness demands evaluating the logical entailment derived strictly from the prompt. Furthermore, the taxonomy introduced by \citet{cattan2025draggedconflictsdetectingaddressing} categorizes these knowledge divergences and reveals that diverse conflict types necessitate targeted behavioral responses. Their benchmark illustrates that forcing a singular resolution often degrades evaluation quality. 

Addressing these factual inconsistencies directly, the architecture developed by \citet{liu2025truthfulragresolvingfactuallevelconflicts} utilizes structured knowledge graphs to filter contradictory elements systematically. By converting unstructured retrieved text into relational triplets, their method isolates the exact locus of the paradox, ensuring that the evaluator maintains a stable anchor when navigating noisy submissions.

\subsection{Probing the Black Box}
Understanding the failure modes of evaluators requires examining their internal mechanistic responses to contradictory information rather than merely observing their final outputs. Researchers have begun deploying representational probes to monitor neural activations during the exact moment a conflict occurs. 

The investigative work by \citet{gao2025probinglatentknowledgeconflict} identifies that knowledge reconciliation occurs predominantly within the middle layers of the neural architecture. Their research uncovers a latent conflict signal embedded in the hidden states, which is capable of predicting whether the model will ultimately prioritize its training weights or the provided context. 

Complementing this foundational discovery, the representational analysis conducted by \citet{zhao2025understanding} details the specific competitive dynamics between attention mechanisms and multilayer perceptrons. They demonstrate that while attention heads attempt to heavily weight the external evidence, the multilayer perceptrons continuously inject the original parametric priors, causing a measurable internal collision. The transparency methods established by \citet{ye2026seeingconflicttransparentknowledge} further elucidate these internal compromises. By mapping the precise attention flow and activation patterns, their approach makes the underlying epistemological struggle completely observable at the token generation level, proving that evaluators do not hallucinate randomly but rather fail through predictable mechanistic pathways.

\subsection{Dynamic Reliance and Control}
The recognition of internal conflict mechanisms enables the development of dynamic strategies for behavioral control, moving beyond passive observation to active intervention. 

The precise intervention techniques explored by \citet{bi2025parametersvscontextfinegrained} demonstrate that manipulating specific network components allows for exact regulation of the reliance on either parametric or contextual knowledge. By adjusting activation thresholds in targeted layers, engineers can force the evaluator to temporarily suppress its pretraining bias when assessing highly specialized or novel documents. Building on these defensive capabilities, the reinforcement methodology presented by \citet{lin2026resistingcontextualinterferencerag} fortifies foundational knowledge to withstand malicious or noisy contextual interference. Their training paradigm ensures that the core reasoning pathways remain robust even when flooded with irrelevant distractors. 

Scaling these controls to complex environments, the adaptive framework designed by \citet{wang2025accommodateknowledgeconflictsretrievalaugmented} facilitates reliable response generation amidst unpredictable practical data conflicts. This system dynamically shifts its confidence weighting based on the detected quality of the retrieved context. Finally, the actionable internal reasoning approach proposed by \citet{huo2025microactmitigatingknowledgeconflict} empowers evaluators to resolve discrepancies actively. Instead of relying on static hidden state adjustments, this method encourages the model to generate explicit intermediate logic steps, marking a crucial transition toward autonomous verification and paving the way for advanced agentic assessment systems.